\section*{Hoofdstuk \ref{ch:g3:continents-and-oceans} --- De globe: continenten en oceanen}

\begin{solution}{exo:g3:continents-and-oceans:1}
Azië, Afrika, Noord-Amerika, Zuid-Amerika, Antarctica, Europa en
Australië. Het grootste is Azië; het kleinste is Australië.
\end{solution}

\begin{solution}{exo:g3:continents-and-oceans:2}
De Stille, de Atlantische, de Indische en de Zuidelijke \omterm{def:g2:water-on-earth:ocean}{Oceaan}, en de
Noordelijke IJszee. De grootste is de Stille \omterm{def:g2:water-on-earth:ocean}{Oceaan}; de kleinste is de
Noordelijke IJszee.
\end{solution}

\begin{solution}{exo:g3:continents-and-oceans:3}
De \omterm{def:g3:continents-and-oceans:poles}{Noordpool} ligt bovenaan op een \omterm{def:g3:continents-and-oceans:globe}{globe}; de \omterm{def:g3:continents-and-oceans:poles}{Zuidpool} ligt onderaan; de
\omterm{def:g3:continents-and-oceans:equator}{evenaar} loopt helemaal om het midden heen.
\end{solution}

\begin{solution}{exo:g3:continents-and-oceans:4}
Europa: noordelijk. Australië: zuidelijk. Antarctica: zuidelijk.
Noord-Amerika: noordelijk.
\end{solution}

\begin{solution}{exo:g3:continents-and-oceans:5}
De Noordelijke IJszee ligt rond de \omterm{def:g3:continents-and-oceans:poles}{Noordpool}. De Zuidelijke \omterm{def:g2:water-on-earth:ocean}{Oceaan} ligt
rond Antarctica.
\end{solution}

\begin{solution}{exo:g3:continents-and-oceans:6}
Van Afrika naar Zuid-Amerika: de Atlantische \omterm{def:g2:water-on-earth:ocean}{Oceaan}. Van Australië naar
het oosten, naar Zuid-Amerika: de Stille \omterm{def:g2:water-on-earth:ocean}{Oceaan}. Van Afrika naar
Australië: de Indische \omterm{def:g2:water-on-earth:ocean}{Oceaan}.
\end{solution}

\begin{solution}{exo:g3:continents-and-oceans:7}
Sam woont bij de \omterm{def:g3:continents-and-oceans:equator}{evenaar}, waar het het hele jaar door warm is. Bo woont bij
een \omterm{def:g3:continents-and-oceans:poles}{pool}, waar het het hele jaar door koud is.
\end{solution}

\begin{solution}{exo:g3:continents-and-oceans:8}
$10 - 7 = 3$: haar vinger komt 3 keer op het land. Dat konden we
verwachten: ongeveer 71 vierkantjes van de 100 van de Aarde zijn water, dus
stopt de vinger ongeveer 7 keer op de 10 op de \omterm{def:g2:water-on-earth:ocean}{zee}. (Met maar 10 keer
draaien kan ze ook 6 of 8 krijgen: de telling verschilt van spel tot spel.)
\end{solution}

\begin{solution}{exo:g3:continents-and-oceans:9}
Om de \omterm{def:g3:continents-and-oceans:world-map}{wereldkaart} te tekenen, werd de \omterm{def:g3:continents-and-oceans:globe}{globe} opengeknipt langs een lijn door
het midden van de Stille \omterm{def:g2:water-on-earth:ocean}{Oceaan} en plat gelegd. De ene helft van de Stille
\omterm{def:g2:water-on-earth:ocean}{Oceaan} komt aan de linkerrand van de \omterm{def:g2:drawing-to-map:map}{kaart} terecht, de andere helft aan de
rechterrand. Op de echte Aarde is er maar één Stille \omterm{def:g2:water-on-earth:ocean}{Oceaan}: de twee randen
van de \omterm{def:g2:drawing-to-map:map}{kaart} zijn dezelfde lijn op de \omterm{def:g3:continents-and-oceans:globe}{globe}.
\end{solution}

\begin{solution}{exo:g3:continents-and-oceans:10}
Tom heeft geen gelijk. Van groot naar klein zijn de \omterm{def:g3:continents-and-oceans:continent}{continenten} Azië,
Afrika, Noord-Amerika, Zuid-Amerika, Antarctica, Europa en Australië:
Antarctica is maar het vijfde, veel kleiner dan Azië. De \omterm{def:g3:continents-and-oceans:world-map}{wereldkaart} houdt
hem voor de gek omdat een platte \omterm{def:g2:drawing-to-map:map}{kaart} het land dicht bij de \omterm{def:g3:continents-and-oceans:poles}{polen} flink
uitrekt. Antarctica, dat rond de \omterm{def:g3:continents-and-oceans:poles}{Zuidpool} ligt, is uitgerekt tot een lange
band over de hele onderkant van de \omterm{def:g2:drawing-to-map:map}{kaart}. Op de \omterm{def:g3:continents-and-oceans:globe}{globe} is het een ronde
landmassa, en veel kleiner dan Azië.
\end{solution}
